\documentclass[aspectratio=169, 10pt]{beamer}
\usetheme{metropolis}

% --- Edinburgh colours ----------------------------------------
\definecolor{UoERed}{HTML}{7A2318}
\definecolor{UoEGold}{HTML}{B8860B}
\definecolor{CodeBG}{HTML}{F7F3ED}
\definecolor{CodeFrame}{HTML}{D5C9B1}

\setbeamercolor{palette primary}{bg=UoERed, fg=white}
\setbeamercolor{frametitle}{bg=UoERed, fg=white}
\setbeamercolor{title separator}{fg=UoEGold}
\setbeamercolor{progress bar}{fg=UoEGold, bg=UoEGold!20}
\setbeamercolor{alerted text}{fg=UoERed}
\setbeamercolor{block title}{bg=UoERed!15, fg=UoERed}
\setbeamercolor{block body}{bg=UoERed!5}

% --- Packages -------------------------------------------------
\usepackage[T1]{fontenc}
\usepackage{lmodern}
\usepackage{amsmath, amssymb, mathtools}
\usepackage{listings}
\usepackage{graphicx, booktabs}
\usepackage{tikz}
\usetikzlibrary{arrows.meta, positioning, calc, decorations.pathreplacing}
\usepackage{hyperref}
\hypersetup{colorlinks=true, linkcolor=UoERed, urlcolor=UoERed!70!black}
\setbeamertemplate{navigation symbols}{}

\lstset{
  language=Python,
  basicstyle=\ttfamily\scriptsize,
  keywordstyle=\color{UoERed}\bfseries,
  commentstyle=\color{gray}\itshape,
  stringstyle=\color{UoEGold},
  backgroundcolor=\color{CodeBG},
  frame=single,
  rulecolor=\color{CodeFrame},
  numbers=none, breaklines=true, showstringspaces=false, tabsize=4,
  xleftmargin=4pt, xrightmargin=4pt, aboveskip=6pt, belowskip=4pt,
}

\AtBeginSection[]{
  \begin{frame}
    \vfill\centering
    \usebeamerfont{title}\insertsectionhead\par
    \vfill
  \end{frame}
}

\title{Week 10 --- Frontiers \& Course Synthesis}
\subtitle{Sequence-Space Methods, Continuous-Time HA \& When to Use Deep Learning}
\author{Dr Juan Zurita}
\institute{Deep Learning for Macroeconomics\\
The University of Edinburgh~$\cdot$~School of Economics}
\date{}

\begin{document}
\maketitle

\begin{frame}{Sequence-Space Jacobians}
\textbf{Auclert et al.\ (2021):} solve HA models by linearising around the steady state in \textbf{sequence space}.

\bigskip
\begin{itemize}
\item Represent the economy as a mapping from shock sequences to outcome sequences
\item Compute Jacobians of this mapping using autodiff
\item Linear impulse responses without simulation
\end{itemize}

\bigskip
Complements DEQNs: SSJ is fast for linearised models; DEQNs handle non-linearities.
\end{frame}

\begin{frame}{Continuous-Time HA Models}
\textbf{Achdou et al.\ (2022):} solve the Kolmogorov Forward Equation (KFE) for the wealth distribution.

\bigskip
$0 = -\partial_a[s(a,z)\,g(a,z)] + \text{(diffusion)} + \text{(jumps in } z\text{)}$

\bigskip
Combined with the HJB for individual decisions:

\bigskip
\begin{itemize}
\item PINNs for the HJB $\Rightarrow$ policy functions
\item Finite differences or PINNs for the KFE $\Rightarrow$ distributions
\item Mean-field games: the two PDEs coupled
\end{itemize}
\end{frame}

\begin{frame}{When to Use Deep Learning}
\begin{center}
\begin{tabular}{ll}
\toprule
Use deep learning when \ldots & Use traditional methods when \ldots \\
\midrule
State dimension $> 4$--$5$ & Low-dimensional ($\leq 3$) \\
Constraints bind occasionally & Linear or log-linear model \\
Non-linear dynamics matter & Perturbation suffices \\
Global solution needed & Local approximation is fine \\
\bottomrule
\end{tabular}
\end{center}

\bigskip
Deep learning is a \textbf{complement} to traditional methods, not a replacement. Use the simplest tool that answers your question.
\end{frame}

\begin{frame}{Course Summary}
\begin{center}
\begin{tabular}{ll}
\toprule
Weeks & Topic \\
\midrule
1--3 & Foundations: networks, training, autodiff \\
4--5 & DEQNs: deterministic and stochastic \\
6 & Constraints: Fischer--Burmeister \\
7 & PINNs: continuous-time models \\
8 & Heterogeneous agents \\
9 & Surrogates and estimation \\
10 & Frontiers \\
\bottomrule
\end{tabular}
\end{center}

\bigskip
From PNM's grids to neural networks --- you now have tools for problems that were infeasible five years ago.
\end{frame}

\begin{frame}{Where to Go Next}
\begin{itemize}
\item \textbf{PNM:} revisit with deeper understanding of what the algorithms do
\item \textbf{Optimisation course:} the mathematical theory behind gradient descent and beyond
\item \textbf{Research:} Scheidegger's GitHub, QuantEcon, Econ-ARK
\end{itemize}

\bigskip
\textbf{References:} Azinovic et al.\ (2022), Scheidegger (2025), Goodfellow et al.\ (2016)

\par\bigskip
\alert{Thank you and good luck!}
\end{frame}

\end{document}
